\documentclass[11pt,a4paper]{article}
\usepackage[utf8]{inputenc}
\usepackage[spanish]{babel}
\usepackage{geometry}
\geometry{margin=2.5cm}

\title{\textbf{Reporte Ejecutivo de Gestión de Riesgos}\\ \large Clínica Privada}
\author{Ignacio De Paz --- LU: 123456}
\date{\today}

\begin{document}
\maketitle

\section{Resumen Ejecutivo}
La gestión de la seguridad de la información en nuestra institución médica demanda un enfoque preventivo frente al incremento de ciberamenazas en el sector salud. Al igual que en la fábula del \textbf{Lobo Feroz y Caperucita Roja}, las amenazas externas acechan continuamente disfrazadas de comunicaciones legítimas (como los ataques de phishing) para engañar al personal y vulnerar el perímetro de la clínica.

Para contrarrestar estos riesgos, debemos adoptar la estrategia de los \textbf{Tres Cerditos}: construir una infraestructura sólida de ladrillo mediante la implementación de defensas en profundidad y capas múltiples de seguridad (autenticación multifactor, segmentación de red y copias de seguridad inmutables), evitando depender de estructuras frágiles que sucumban ante el primer embate.

\section{Principales Riesgos Identificados}
\begin{enumerate}
    \item \textbf{Compromiso de credenciales de HCD por Phishing (Nivel 20 - Crítico):} Vulnerabilidad en el acceso a historias clínicas digitales.
    \item \textbf{Infección por Ransomware en Servidor de Imágenes Médicas (Nivel 20 - Crítico):} Amenaza directa a la continuidad operativa del sistema PACS.
    \item \textbf{Interrupción del Suministro Eléctrico en Data Center (Nivel 12 - Alto):} Riesgo sobre la disponibilidad de la infraestructura central.
\end{enumerate}

\section{Recomendaciones Prioritarias y Plan de Acción}
\begin{itemize}
    \item \textbf{Despliegue de Autenticación Multi-Factor (MFA):} Implementación obligatoria para el personal médico.
    \item \textbf{Backup Inmutable y Segregación de Red:} Aislamiento de sistemas críticos contra ataques de ransomware.
    \item \textbf{Adecuación Física de las Instalaciones:} Como medida prioritaria de seguridad física, se establece estrictamente que los \textbf{centros de datos no deben ubicarse junto a cocinas ni debajo de piletas} o cañerías de agua, previniendo riesgos críticos por filtraciones, humedad o incendios.
\end{itemize}

\end{document}
